\section{Fourier estimates and the range \texorpdfstring{\(s\ge d\)}{s >= d}}
\label{sec:large-s}

Identify the torus \(\T^d\) with the half-open unit cube and use
its volume-one measure. Periodize the zero extension of \(f\).
It agrees with \(f\) almost everywhere on the cube and has
\(\|f\|_2\le L\) and mean one. We use \(f\) also for the periodized
function from here on. Constants \(c,C\) may change between
occurrences and depend on the fixed class parameters. Further
fixed parameters in the construction are chosen in terms of those
parameters.

For \(|u|\le\sqrt d\), the supports of the terms of the
whole-space difference lie in a fixed box. Hölder's inequality in
that box bounds its \(L^2\) norm by \(C|u|^s\). The periodization
of the difference is the torus difference. Only a bounded number
of terms can overlap at each point in this periodization, so
Cauchy--Schwarz gives \(\|\Delta_u^r f\|_{2,\T^d}\le C|u|^s\).
On a Fourier mode \(k\in\mathbb Z^d\), the centered difference
has multiplier \((2i\sin(\pi k\cdot u))^r\). For \(M\ge1\), average
the squared norm over \(u\in[-M^{-1},M^{-1}]^d\). At any
frequency \(|k|_\infty>M\), one coordinate of \(k\cdot u\)
ranges over an interval longer than two as that coordinate of \(u\)
varies. The average of \(|2\sin(\pi k\cdot u)|^{2r}\) is
bounded below by a positive constant, by retaining the complete
periods in that interval. Parseval and Tonelli give
\begin{equation}\label{eq:tail}
   \sum_{|k|_\infty>M} |\widehat f(k)|^2\le C M^{-2s}.
\end{equation}
All estimates use the original difference order, including \(r=s\)
at integer \(s\). In the real Fourier shell
\(2^{j-1}<|k|_\infty\le2^j\) the dimension is \(O(2^{jd})\)
and the \(L^2\) norm is \(O(2^{-js})\). Use the shell
\(0<|k|_\infty\le1\) when \(j=0\).

To construct the mixture, fix \(M=2^J\), \(\zeta=M^{-s}\).
In shell \(j\le J\) choose a deterministic
\(\zeta 2^{-(J-j)}\)-net of its coefficient ball of radius
\(C2^{-js}\). The logarithm of the covering cardinality of a
radius-\(R\) ball in dimension \(m\) at resolution \(\rho\)
is at most \(m\log(1+2R/\rho)\), by comparing disjoint
radius-\(\rho/2\) balls centered at a maximal separated set.
Sum one choice from each shell and adjoin the constant one to
obtain the family \(\mathcal G_M\). Its members have uniformly
bounded \(L^2\) norms. By orthogonality and \eqref{eq:tail},
each \(f\) has a representative \(g\in\mathcal G_M\)
with \(\|g-f\|_2\le C\zeta\). The covering count is
\[
 \log|\mathcal G_M|\le C\sum_{j=0}^J2^{jd}(1+J-j)\le C M^d .
\]
For each \(g\), put \(q_g=(g_++\zeta)/Z_g\),
where \(Z_g=\int(g_++\zeta)\) is in \([1,C]\). For a representative
\(g\) of \(f\), truncation reduces the distance to \(f\),
\(|Z_g-1|\le \|g_++\zeta-f\|_1\le C\zeta\), and hence
\[
 \|q_g-f\|_2\le C\zeta,\qquad
 \chi^2(f,q_g):=\int (f-q_g)^2/q_g\le C\zeta .
\]
The floor \(q_g\ge c\zeta\) applies to the likelihoods used
in the estimator.

We use the finite-net likelihood comparison of Yang and Barron
\cite[Section~2, equation (2)]{YangBarron} as a pointwise
comparison and bound both signs in \(L^2\). Let
\(Q_n=|\mathcal G_M|^{-1}\sum_{g\in\mathcal G_M} q_g^{\otimes n}\)
and use \(\widehat H=-n^{-1}\log Q_n(X_1,\ldots,X_n)\).
It is finite and measurable on the product cube. If
\(B=n^{-1}\log(f^{\otimes n}/Q_n)(X_1,\ldots,X_n)\), then
\(\E_f e^{-nB}\le1\), by integrating the density \(Q_n\)
over the set where \(f^{\otimes n}>0\). Thus
\(\PP_f(B<-t)\le e^{-nt}\) and \(\|B_-\|_2\le \sqrt{2}/n\).
For the representative \(q=q_g\), the corresponding term
in the mixture bounds \(B\) above by
\[
 n^{-1}\log|\mathcal G_M|
 +n^{-1}\sum_{i=1}^n\log(f(X_i)/q(X_i)).
\]
The last average has mean in \([0,\chi^2(f,q)]\) by Jensen
and \(\log z\le z-1\). Its variance is at most \(C\zeta/n\)
since \(z(\log z)^2\le C(z-1)^2\) for \(z\ge0\), by
continuity at \(0,1\) and the bound at infinity. Consequently
\(\|B_+\|_2\le CM^d/n+C\zeta+C\sqrt{\zeta/n}\).
The empirical mean of \(-\log f(X_i)\) has mean \(H(f)\)
and standard deviation at most \(Cn^{-1/2}\), since
\(z(\log z)^2\le C(1+z^2)\). Adding \(B\) to that mean yields
\[
 \|\widehat H-H(f)\|_{L^2(\PP_f)}
 \le C\bigl(n^{-1/2}+M^{-s}+M^d/n\bigr).
\]
Take \(M\asymp n^{1/(s+d)}\) dyadic to get \(Cn^{-1/2}\)
for every fixed \(s\ge d\), including \(s=d\).
